\scanpage{56}
\mspage{27}

. Le fait que $\lambda$ centralise $G_0$ implique que l'on a, pour $u\in G_0$

\[
 \{s',t'\}\in A_\lambda\Longleftrightarrow\{us',ut'\}\in A_\lambda
 \tag*{(142)}
\]

puisque les deux termes de l'équivalence s'écrivent resp.$^{\mathrm t}$

\[
 \{s',\lambda^{-1}t'\}\in A_0\qquad\text{et}\qquad
 \{us',\underbrace{\lambda^{-1}ut^{\unsure{+|\prime}}}_{u\lambda^{-1}t'}\}\in A_0,
\]

relations dont l'équivalence signifie justement que $u\in G_0$. Donc on trouve encore

\[
 \deleted{\unsure{G_\lambda}}\quad\underbrace{\operatorname{Aut}(S,A_\lambda)}_{G_\lambda}=G_0
 \tag*{(143)}
\]

\uline{Définition} Soient $S$ un ensemble, $A_0\subset\unsure{\mathscr P|\mathcal P|\mathfrak P}_2(S)$, $A_1\in\mathscr P_2(S)$ deux structures de graphes \uline{stricts} \inserted{(i.e. une arête connue quand on connaît l'ens. de ses \unsure{extrémités|extrêmités}, et celles-ci sont distinctes)} sur $S$, \inserted{admettant} comme ens. de ses sommets, \inserted{$G_0$ et $G_1$ les groupes des autom. \unsure{correspts.},} On dit que ces structures sont \uline{équivalentes}, s'il existe un $\lambda\in\unsure{\mathfrak S|\mathscr S}_S$, \unsure{telle|tel} que \deleted{l'on} l'on ait

a) $\lambda$ centralise $G_0$

b) \deleted{\unsure{$s,t$}} \deleted{\unsure{Sois}} $s,t\in S$, alors

\[
 \deleted{\unsure{\{\lambda^{-1}s',}}\quad
 \{s,t\}\in A_0\Longleftrightarrow\{s,\lambda t\}\in A_1
\]

% slanted

\margin{On dit alors que $\lambda$ est un \uline{transporteur} de $A_0$ vers $A_1$}

\uline{Remarque} L'argument précédent montre que l'on a alors

\[
 G_0=G_1,
\]

% archive sheet number: 55
